% ---------------------------------------------------------------------
%  Chapter 11 : Caching --- the enumeration and expand caches
% ---------------------------------------------------------------------
\chapter{Caching: The Enumeration and Expand Caches}\label{ch:caching}

Two of the steps on the road from a model spec to a number are
\emph{brutally expensive} and \emph{highly reusable}, and this chapter is
about how \Daedalus{} avoids paying for them twice. The first is expanding
the symbolic action: a multivariate Taylor expansion that, for a four-field,
two-population compartmental theory at \code{taylor\_order=4}, runs
\emph{roughly ninety minutes single-threaded} inside SageMath. The second is
enumerating Feynman diagrams: a four-stage combinatorial pipeline whose cost
grows fast with the external-leg count $k$ and the loop order $\ell$.

The key observation that makes caching possible --- and the organising idea
of the whole chapter --- is that \emph{neither result depends on the numerical
parameter values a user later plugs in}. A higher value of the friction
$\mu$, a different noise strength $D$: these change the numbers that come out
the far end of the pipeline, but they do not change \emph{which} diagrams
exist or \emph{what} the expanded action looks like symbolically. Those
artefacts depend only on the \emph{structure} of the theory --- its field
list, its interaction vertices, the requested $k$ and $\ell$, and which fields
sit on the external legs. So both are computed once, written to disk keyed by
that structure, and reloaded across notebook restarts and across runs that
vary only the numbers.

This chapter documents the caching layer for those two artefacts. It sits
squarely between the model definition and the numerical integration, and it is
built on four files:

\begin{description}
  \item[\file{msrjd/core/cache.py}] the low-level, stage-keyed \code{.sobj}
        store, \code{PipelineCache}, that both higher-level caches build on.
  \item[\file{pipeline/\_expand\_cache.py}] the disk cache for
        \code{FieldTheory.expand()} results --- the expanded action.
  \item[\file{pipeline/\_diagrams.py}] the disk cache wrapping the four-stage
        diagram pipeline --- \code{enumerate\_unique\_diagrams}.
  \item[\file{pipeline/\_precompute.py}] the cheap one-time structural pass,
        \code{precompute}, that warms the durable caches up front.
\end{description}

By the end you will understand exactly which file lands where on disk, why a
cached order-6 expansion satisfies an order-2 request for free, why the
enumeration cache embeds a version string that you \emph{must} bump when you
change the diagram semantics, and which single corner of this subsystem is the
most fragile.

\section{Why cache at all}

\Daedalus{} turns a model spec --- a Python \code{dict} describing an \msrjd{}
field theory --- into the connected cumulants of the fields, computed
perturbatively from Feynman diagrams. Two stages on that road dominate the
cost and are perfectly reusable.

\paragraph{Expanding the action.}
To do perturbation theory you need the action written as a power series in the
\emph{fluctuations} about the mean-field saddle. \Daedalus{} hands the
symbolic \msrjd{} action $S[\tilde\varphi,\varphi]$ to SageMath and runs a
multivariate Taylor expansion to a chosen order. The docstring of
\file{pipeline/\_expand\_cache.py:6--9} is blunt about the cost:

\begin{quote}\itshape
``\code{FieldTheory.expand()} is the dominant cost for non-trivial theories
--- for a 4-field $\times$ 2-population compartmental Bernoulli theory at
\code{taylor\_order=4} it runs $\sim$90 minutes single-threaded inside Sage's
multivariate \code{taylor()}.''
\end{quote}

\paragraph{Enumerating diagrams.}
For each loop order $\ell = 0, 1, \dots, \texttt{max\_ell}$ the pipeline
enumerates every topologically distinct, type-consistent, causal,
\emph{non-isomorphic} Feynman diagram with $k$ external legs. This is a
four-stage combinatorial pipeline --- prediagrams, then typing, then a
causality filter, then deduplication --- whose cost grows fast with $k$ and
$\ell$.

\paragraph{The separation of structure from numbers.}
The thing that makes both cacheable is that they are functions of the theory's
\emph{structure} alone. Concretely, the two artefacts depend on the field
list, the interaction vertices, the requested $k$ and $\ell$, and the
external-field labels --- and on \emph{nothing} numerical. So \Daedalus{}
caches both on disk, keyed by structure, and reloads them across notebook
restarts and across parameter sweeps. This is the single most important fact
in the chapter: \emph{the cache key is structural, never numerical}.

The data path looks like this. \code{TheoryBuilder.build()} produces the model
\code{dict}; \code{precompute(model)} runs a one-time structural pass that
warms the durable caches; and \code{compute\_cumulants(\dots)} consumes them
on its way to the integrator:

\begin{lstlisting}[style=console]
TheoryBuilder.build()  ->  model dict
        |
        v
   precompute(model)  -- one-time structural pass --+
        |   (expand@order2 + propagator + MF check) |  writes
        v                                            v
  compute_cumulants(...)                  saved_theories/<slug>/
        |                                    +-- propagator.sobj
        +-[1] FieldTheory.expand            +-- expand_taylor2.sobj
        |     <- EXPAND CACHE               +-- expand_taylor4.sobj
        |       (pipeline/_expand_cache.py) +-- unique_typed_..._k<k>_l<l>_taylor<N>.sobj
        +-[5] enumerate_unique_diagrams     +-- manifest.json
        |     <- ENUMERATION CACHE
        |       (pipeline/_diagrams.py)
        v
   diagrams -> integration -> cumulants
\end{lstlisting}

The two caches have deliberately \emph{different} keys, and that difference is
the central design decision:

\begin{itemize}
  \item The \term{expand cache} is keyed by \code{taylor\_order} \emph{only}
        --- one file per order. A higher-order expansion is a strict superset
        of a lower one (we prove this below), so a high-order cache serves any
        lower-order request for free.
  \item The \term{enumeration cache} is keyed by the finer tuple
        $(\text{model}, \texttt{taylor\_order}, k, \ell,
        \text{external-fields})$, because the diagram \emph{set} does depend on
        which interaction vertices are in scope, and that in turn depends on
        \code{taylor\_order}.
\end{itemize}

Both caches share one physical store, \code{PipelineCache}, which writes Sage
\code{.sobj} files into a single per-theory directory and maintains a
human-readable \code{manifest.json} index. (As we will see, the expand cache
actually bypasses \code{PipelineCache} and calls Sage's \code{save}/\code{load}
directly --- but it writes into the \emph{same} directory.)

\section{The on-disk layout, and \texttt{.sobj} files from scratch}

Everything for one theory lives in a single directory under
\file{saved\_theories/}. The directory name is the theory's \term{slug}: a
filesystem-safe version of \code{model['name']}. All three caches compute the
same slug the same way, which is why they share one directory. The rule is
(\file{pipeline/\_expand\_cache.py:104--106}):

\begin{lstlisting}
def _slug(model: dict) -> str:
    """Filesystem-safe slug for the cache directory."""
    return re.sub(r'[^A-Za-z0-9]+', '_', model['name']).strip('_').lower()
\end{lstlisting}

The regular expression \code{[\^{}A-Za-z0-9]+} matches any run of characters
that are \emph{not} a letter or digit, and \code{re.sub} replaces each such run
with a single underscore; \code{.strip('\_')} trims leading and trailing
underscores, and \code{.lower()} normalises case. So a model named
\code{'1d KPZ (per-leg gradient vertex)'} becomes the directory
\file{1d\_kpz\_per\_leg\_gradient\_vertex}. A live such directory looks like:

\begin{lstlisting}[style=console]
saved_theories/1d_kpz_per_leg_gradient_vertex/
    expand_taylor2.sobj        propagator.sobj
    expand_taylor4.sobj        manifest.json
    unique_typed_mult_v3_<ext>_k<k>_l<l>_taylor<N>.sobj   (one per (k, l, ext, order))
\end{lstlisting}

Read the filenames left to right. \file{expand\_taylor2.sobj} and
\file{expand\_taylor4.sobj} are two different orders of the same expanded
action --- siblings in one directory. \file{propagator.sobj} is the symbolic
propagator (built by a different subsystem, but landing in the same place
because it shares the slug). \file{manifest.json} is the human-readable index
of the \code{PipelineCache} entries. And the long
\file{unique\_typed\_mult\_v3\_\dots} files are the per-$(k,\ell)$ enumeration
results; each filename carries the external-field tag, the version stamp, the
$k$ and $\ell$, and the taylor order. We dissect every piece of that name in
Section~\ref{sec:enumcache}.

\begin{defn}
A \term{\code{.sobj}} (``Sage object'') file is SageMath's serialization
format: a gzip-compressed pickle written by \code{sage.all.save(obj, path)}
and read back by \code{sage.all.load(path)}. It is the on-disk format for
\emph{every} cached artefact in this subsystem.
\end{defn}

\subsection{Why Sage's serializer, and not plain \texttt{pickle}}

The objects this subsystem writes to disk are \emph{Sage objects}: the
symbolic action, the polynomial-ring elements stored in the expanded action,
the propagator matrix, and the diagram objects (which wrap Sage graphs). Plain
Python \code{pickle} does not round-trip these reliably, because Sage objects
carry references to \emph{parent structures} --- rings, fields --- whose
identity matters. After an unpickle, the reconstructed parent is a
\emph{different} Python object than a freshly built one, even when the two have
identical generators, and arithmetic that mixes elements from the two parents
can break. SageMath's own \code{save}/\code{load} understand this and handle it
correctly, so \Daedalus{} uses them everywhere it touches disk.

\begin{note}
SageMath is a large open-source mathematics system built on top of Python. It
bundles symbolic algebra (its symbolic ring \code{SR}), exact and
arbitrary-precision arithmetic, polynomial rings, and graph theory under one
Python API. When you import from \code{sage.all} you get its global namespace,
including the two object-serialization helpers used throughout this chapter:
\code{save} and \code{load}. If you have never used Sage, the only two facts
you need for this chapter are: (i) \code{sage.all.save(obj, path)} pickles
\code{obj} to \code{path.sobj}, and (ii) it \emph{appends} the \code{.sobj}
extension if you do not supply one --- which is why every call here strips it
off first.
\end{note}

There is one subtlety in how \code{save} is called that recurs in every file.
Because Sage's \code{save} appends \code{.sobj}, the caller passes a path with
the extension already \emph{removed}, so that auto-append produces exactly one
\code{.sobj} and not \code{....sobj.sobj}:

\begin{lstlisting}
sage_save(data, path.removesuffix('.sobj'))  # sage_save adds .sobj
\end{lstlisting}

That exact line appears at \file{msrjd/core/cache.py:108},
\file{pipeline/\_expand\_cache.py:345}, and inside \code{save\_expand}. The
matching \code{sage\_load(path)} accepts \emph{either} form (with or without
the extension), so the load side passes the full \code{.sobj} path.

\section{\texttt{PipelineCache}: the low-level store}

\file{msrjd/core/cache.py} is the foundation both higher-level caches sit on.
It is a stage-based disk cache: one directory per model, with files keyed by
the triple \code{(stage\_name, k, loop\_order)}. The class is small enough to
walk method by method.

\subsection{Construction and the stage key}

\code{PipelineCache.\_\_init\_\_} (\file{cache.py:57}) takes a \code{root}
directory and stores \code{self.root = os.path.expanduser(root)} --- the
\code{expanduser} call turns a leading \code{\~{}} into the user's home
directory. It does \emph{not} create the directory yet; that happens lazily on
the first \code{save}.

A cache file's name is built from the stage name plus the optional $k$ and
$\ell$. Two helpers do this (\file{cache.py:62--83}):

\begin{lstlisting}
@staticmethod
def _to_int(val):
    """Coerce SageMath Integer / numpy int to plain Python int."""
    if val is None:
        return None
    return int(val)

@classmethod
def _stage_key(cls, stage, k=None, loop_order=None):
    """Build a filename stem from stage name and optional (k, l)."""
    k = cls._to_int(k)
    loop_order = cls._to_int(loop_order)
    parts = [stage]
    if k is not None:
        parts.append(f'k{k}')
    if loop_order is not None:
        parts.append(f'l{loop_order}')
    return '_'.join(parts)
\end{lstlisting}

\code{\_to\_int} earns its place because $k$ and $\ell$ frequently arrive as
SageMath \code{Integer} objects, not plain Python \code{int}s. Coercing them
keeps filenames and manifest JSON clean --- you get \code{k2}, not the
\code{repr} of a Sage \code{Integer}. So \code{('unique\_typed', 2, 1)} becomes
the stem \code{'unique\_typed\_k2\_l1'}, and \code{\_sobj\_path}
(\file{cache.py:81}) joins that to the root with a \code{.sobj} suffix.

\subsection{Save, load, exists, and the memoise helper}

The core API is four methods. \code{exists} (\file{cache.py:87}) is a one-liner
--- \code{os.path.isfile(self.\_sobj\_path(\dots))} --- and it is the
\emph{filesystem} that is consulted, not the manifest. \code{save}
(\file{cache.py:91}) creates the directory if needed, writes via Sage, and
then updates the manifest:

\begin{lstlisting}
def save(self, stage, data, k=None, loop_order=None):
    os.makedirs(self.root, exist_ok=True)
    path = self._sobj_path(stage, k, loop_order)
    sage_save(data, path.removesuffix('.sobj'))  # sage_save adds .sobj
    self._update_manifest(stage, k, loop_order)
\end{lstlisting}

The \code{exist\_ok=True} flag on \code{os.makedirs} means ``do not raise if
the directory already exists'' --- the idiomatic way to ensure a directory is
present without a prior existence check. \code{load} (\file{cache.py:111})
reads the file back through \code{sage\_load}, and \emph{raises}
\code{FileNotFoundError} with a descriptive message if the file is absent. That
exception is the documented contract, but --- as we will see --- the
enumeration cache calls \code{exists} first and wraps \code{load} in a
\code{try}/\code{except}, so it never actually relies on the raise.

\code{get\_or\_compute} (\file{cache.py:125}) is the classic memoise pattern:

\begin{lstlisting}
def get_or_compute(self, stage, compute_fn, k=None, loop_order=None):
    if self.exists(stage, k, loop_order):
        return self.load(stage, k, loop_order)
    data = compute_fn()
    self.save(stage, data, k, loop_order)
    return data
\end{lstlisting}

\begin{note}
The enumeration cache does \emph{not} use \code{get\_or\_compute}. It does its
own manual \code{exists}/\code{load}/\code{save} dance so that it can print
per-stage progress messages and degrade gracefully on a load failure. The
memoise helper is here for callers who want the simple path; the production
diagram cache wants the verbose, fault-tolerant one.
\end{note}

\subsection{The manifest --- a best-effort index}

\code{PipelineCache} maintains a \file{manifest.json} alongside the
\code{.sobj} files: a human-readable index of what has been cached. Its shape
is \code{\{'entries': [\{key, stage, k, loop\_order, saved\_at\}], 'created':
iso, 'updated': iso\}}. \code{\_update\_manifest} (\file{cache.py:168}) loads
the current manifest, removes any stale entry with the same key, appends a
fresh one with an ISO-8601 timestamp, and writes it back with
\code{json.dump(\dots, indent=2)} for readability.

The crucial design point is that the manifest is \emph{not} correctness-bearing.
\code{\_load\_manifest} (\file{cache.py:153}) is deliberately robust: if the
file is missing \emph{or corrupt} it silently returns a fresh empty structure
rather than raising:

\begin{lstlisting}
try:
    with open(mp) as f:
        manifest = json.load(f)
    if not isinstance(manifest, dict):
        raise ValueError
    return manifest
except (json.JSONDecodeError, ValueError):
    # Corrupt manifest -- start fresh.
    return {'entries': [], 'created': datetime.now().isoformat()}
\end{lstlisting}

The \code{.sobj} files are the source of truth; the manifest is a convenience.
Because \code{exists} checks the filesystem and not the manifest, a missing or
truncated manifest entry can never cause a spurious cache miss.

Finally, \code{clear} (\file{cache.py:194}) removes cached data: with all three
arguments \code{None} it \code{shutil.rmtree}s the whole directory; otherwise
it deletes just the one matching \code{.sobj} and prunes its manifest entry.
And \code{\_\_repr\_\_} reports \code{PipelineCache('<root>', <N> entries)}.

\section{The expand cache I: bigrades and the superset theorem}

We now turn to the larger of the two caches, the one that stores the
ninety-minute artefact. To understand what it stores you first need to
understand the object \code{FieldTheory.expand()} produces: the
bigrade-classified action.

\subsection{The \msrjd{} action and its Taylor expansion}

\Daedalus{} works in the Martin--Siggia--Rose--Janssen--De~Dominicis
representation of a stochastic dynamical system. A Langevin equation
$\partial_t\varphi = F[\varphi] + \text{noise}$ becomes a path integral with an
action over a \emph{physical} field $\varphi$ and a \emph{response} (or
``tilde'') field $\tilde\varphi$:
\begin{equation}
  S[\tilde\varphi,\varphi]
  = \int \dd t\,
    \Big\{
      \tilde\varphi\,(\Dt\varphi - F[\varphi])
      - \tfrac12\,\tilde\varphi\,B\,\tilde\varphi + \cdots
    \Big\}.
  \label{eq:msrjd-action}
\end{equation}
Cumulants are read off from a perturbative expansion of $\ee^{-S}$ around the
mean-field (saddle) point. To do that expansion you need the action written in
powers of the fluctuations about the saddle. Write
$\varphi = \varphi^{*} + \delta\varphi$ and
$\tilde\varphi = \tilde\varphi^{*} + \delta\tilde\varphi$, and Taylor-expand $S$
as a multivariate series in $(\delta\tilde\varphi, \delta\varphi)$.

The central move is then to \emph{classify each monomial by its bigrade}.

\begin{defn}
The \term{bigrade} of an action monomial is the pair
\[
  (n_{\tilde{}},\, n_{\mathrm{phys}})
  = (\#\,\text{response-field factors},\ \#\,\text{physical-field factors}).
\]
Its total degree is $n_{\tilde{}} + n_{\mathrm{phys}}$.
\end{defn}

The expanded action is stored as a dictionary, \code{ft.\_by\_tp}, that maps
each bigrade to the \emph{sum} of all action monomials of that bigrade:
\[
  \texttt{\_by\_tp}[(n_{\tilde{}},\, n_{\mathrm{phys}})]
  \;=\; \sum (\text{action monomials of that bigrade}),
\]
each value being a live Sage polynomial-ring element. A Taylor truncation at
order $N$ keeps exactly the bigrades with
$n_{\tilde{}} + n_{\mathrm{phys}} \le N$.

The bigrades carry direct physical meaning, and knowing it tells you why
order~2 is special (Section~\ref{sec:precompute}):

\begin{itemize}
  \item \textbf{$(0,0)$, $(1,0)$, $(0,1)$} --- the \term{mean-field / saddle
        sectors}. These are order zero in the fluctuations: the classical
        equation of motion. The saddle condition is precisely that they vanish
        at $(\varphi^{*}, \tilde\varphi^{*})$, which is what
        \code{sanity\_check} verifies.
  \item \textbf{$(1,1)$} --- the \term{propagator kernel}, the bilinear sector.
        The free (Gaussian) propagator $\Gret$ is the matrix inverse of the
        $(1,1)$ coefficients. The propagator builder reads exactly this sector.
  \item \textbf{Total degree $\ge 3$} --- the \term{interaction vertices}. A
        monomial of bigrade $(a,b)$ with $a+b=d$ is a degree-$d$ vertex with
        $a$ response legs and $b$ physical legs. These are the Feynman rules
        the diagram pipeline consumes.
\end{itemize}

\subsection{The superset theorem and why downgrade is exact}

Here is the mathematical fact that justifies keying the expand cache by
\code{taylor\_order} alone. The docstring states it at
\file{pipeline/\_expand\_cache.py:14--18}:

\begin{quote}\itshape
``\code{\_by\_tp} at \code{taylor\_order=N} is a STRICT SUPERSET of
\code{\_by\_tp} at any \code{taylor\_order=M} with \code{M <= N}: the bigrade
entries at total degree \code{<= M} are byte-identical (the Taylor truncation
only adds higher-order entries, never modifies lower ones).''
\end{quote}

Stated as a theorem:

\begin{defn}
\textbf{(Superset theorem.)} For $M \le N$,
\[
  \texttt{by\_tp}@N \;\supseteq\; \texttt{by\_tp}@M,
  \qquad
  \texttt{by\_tp}@N\big|_{\{(a,b)\,:\,a+b\le M\}}
  \;=\; \texttt{by\_tp}@M .
\]
\end{defn}

The reason is elementary: Taylor truncation at a higher order adds
higher-degree terms but cannot alter the lower-degree ones, because the
order-$d$ Taylor coefficient is fixed by the $d$-th derivative at the saddle and
is independent of where you truncate. Each bigrade entry is therefore
\emph{byte-identical} between the two expansions.

The practical consequence is the \term{downgrade}: to satisfy an order-$M$
request from a cached order-$N$ file, you simply \emph{discard} every bigrade of
total degree above $M$. This is exact --- it is not a re-expansion and not an
approximation. The implementation is one comprehension
(\file{pipeline/\_expand\_cache.py:181--190}):

\begin{lstlisting}
def downgrade_by_tp_dict(by_tp_dict: dict, target_order: int) -> dict:
    """Keep only bigrades whose total degree (a + b) is <= target_order."""
    return {key: poly_dict
            for key, poly_dict in by_tp_dict.items()
            if sum(key) <= target_order}
\end{lstlisting}

Note \code{sum(key)}: \code{key} is the bigrade tuple $(a,b)$, so \code{sum}
gives the total degree $a+b$. This is why a cached order-6 file fully satisfies
an order-2 request, and it is why the cache picks the \emph{smallest} cached
order that suffices: every bigrade kept must later be coerced back into the
freshly built ring, so loading the minimal sufficient file is the cheapest. The
chooser is \code{find\_best\_cached\_order}
(\file{pipeline/\_expand\_cache.py:136--146}):

\begin{lstlisting}
def find_best_cached_order(model, target_order, cache_dir_root='saved_theories'):
    """Return the smallest cached order >= target_order, or None."""
    cached = list_cached_orders(model, cache_dir_root)
    candidates = [o for o in cached if o >= target_order]
    return min(candidates) if candidates else None
\end{lstlisting}

\code{list\_cached\_orders} (\file{pipeline/\_expand\_cache.py:121--133}) finds
all available orders by listing the directory and matching each filename
against \code{re.fullmatch(r'expand\_taylor(\textbackslash d+)\textbackslash
.sobj', fname)} --- the parenthesised \code{(\textbackslash d+)} captures the
integer order, which \code{int(m.group(1))} extracts. It returns \code{[]} if
the directory does not exist.

\begin{gotcha}
Downgrade is exact \emph{only because of} the superset theorem, which in turn
relies on Taylor truncation never modifying lower-order entries. If a future
change made the expansion non-truncating --- a resummation, say --- then
\code{downgrade\_by\_tp\_dict} would silently be \emph{wrong}: it merely drops
high-degree bigrades; it does not recompute the low-degree ones.
\end{gotcha}

\section{The expand cache II: dict-form serialization}

We said Sage objects carry parent-identity baggage across pickle. Polynomial
objects are the worst offenders: an unpickled polynomial ring is a different
Python object than the freshly built one even with identical generators, so a
polynomial unpickled against the old ring will not cleanly combine with
elements of the new ring. The expand cache sidesteps this entirely by
\emph{not storing polynomial objects at all}. Instead it stores their
\term{dict form}.

\begin{defn}
The \term{dict form} of a polynomial-ring element is what Sage's
\code{poly.dict()} method returns: a mapping
\code{\{exponent\_tuple: coefficient\}}. The exponents are plain integers and
the coefficients are \code{SR} (symbolic-ring) elements, both of which survive
pickle and unpickle cleanly.
\end{defn}

The conversions are a handful of one-liners
(\file{pipeline/\_expand\_cache.py:152--175}):

\begin{lstlisting}
def _poly_to_dict_form(poly) -> dict:
    """Returns {exp_tuple: SR_coeff}."""
    return dict(poly.dict())

def _dict_form_to_poly(R, exp_to_coeff: dict):
    """Build a polynomial in ring R from its exp_tuple -> coeff dict."""
    return R(exp_to_coeff) if exp_to_coeff else R.zero()
\end{lstlisting}

The reconstruction calls \code{R(exp\_to\_coeff)} --- the ring's \emph{call
operator}, which builds a polynomial from such a dict. The all-important detail
is that \code{R} here is the \emph{freshly built} ring on load, so the
reconstructed polynomial lives in the live ring and the parent-identity problem
never arises. The whole \code{\_by\_tp} dictionary is mapped through these in
both directions by \code{\_by\_tp\_to\_dict\_form} (save side) and
\code{\_by\_tp\_from\_dict\_form(R, \dots)} (load side, with the rehydrated
ring).

\begin{note}
Because the bundle's values are stored as plain Python dicts of integer
exponents and \code{SR} coefficients, most of the bundle is \emph{pickle-trivial}
--- the only piece where Sage's serializer truly earns its keep is the \code{SR}
coefficients themselves. The dict form thus does most of the
parent-identity work for free, and leaves only the symbolic coefficients to
Sage.
\end{note}

\subsection{The bundle written to disk}

\code{save\_expand} (\file{pipeline/\_expand\_cache.py:322--348}) assembles the
on-disk bundle. It is a plain Python dict:

\begin{lstlisting}
bundle = {
    'by_tp':           _by_tp_to_dict_form(ft._by_tp),
    'S_raw_dict':      (_poly_to_dict_form(ft._S_raw)
                        if ft._S_raw is not None else {}),
    'mf_sector_raw':   _by_tp_to_dict_form(
                          getattr(ft, '_mf_sector_raw', {}) or {}),
    'ring_var_names':  _get_ring_var_names(ft),
    'taylor_order':    int(target),
    'n_tilde':         int(ft._n_tilde),
    'vertex_signature': vertex_form_factor_signature(ft._ns),
    'cache_version':   2,
}
sage_save(bundle, path.removesuffix('.sobj'))
\end{lstlisting}

The fields are summarised in Table~\ref{tab:bundle}.

\begin{table}[t]
\centering\small
\begin{tabular}{@{}lll@{}}
\toprule
key & type & meaning \\
\midrule
\code{'by\_tp'}            & \code{\{(a,b): \{exp: coeff\}\}} & the bigrade-classified action, dict form \\
\code{'S\_raw\_dict'}      & \code{\{exp: coeff\}}            & dict form of the full polynomial (written, \emph{not read on load}) \\
\code{'mf\_sector\_raw'}   & \code{\{(a,b): \{exp: coeff\}\}} & the pre-zero $(0,0)/(1,0)/(0,1)$ sectors (diagnostic) \\
\code{'ring\_var\_names'}  & \code{list[str]}                 & generator names; structural sanity check \\
\code{'taylor\_order'}     & \code{int}                       & the order this was computed at (may exceed the request) \\
\code{'n\_tilde'}          & \code{int}                       & number of response variables; MF-sector filter check \\
\code{'vertex\_signature'} & \code{str} or \code{None}        & operator-IR form-factor fingerprint \\
\code{'cache\_version'}    & \code{int} (=2)                  & bundle-format version \\
\bottomrule
\end{tabular}
\caption{The expand bundle stored in \file{expand\_taylor<N>.sobj}, produced by
\code{save\_expand} and consumed by \code{load\_expand}.}
\label{tab:bundle}
\end{table}

\begin{gotcha}
\code{'S\_raw\_dict'} is written by \code{save\_expand} but \emph{never read by}
\code{load\_expand}. As we will see, the load path reconstructs \code{\_S\_raw}
by summing the (possibly filtered) \code{by\_tp} polynomials instead, which is
the correct value for the requested order. The stored field is therefore
dead-on-load --- harmless, but redundant.
\end{gotcha}

\section{The expand cache III: hidden state and its reconstruction}
\label{sec:hidden-state}

This is the single most fragile corner of the subsystem, and it is worth slowing
down for. The problem is stated plainly in the module docstring
(\file{pipeline/\_expand\_cache.py:67--73}): the on-disk slug is only
\code{model['name']} plus the taylor order, and that key \emph{under-determines}
the theory. Two pieces of state that a fresh \code{expand()} produces as
\emph{side effects} are not captured by \code{\_by\_tp}, hence not by the
bundle, and must be reconstructed when the cache is loaded.

\subsection{The operator-IR form-factor table}

Spatial theories with derivative interaction vertices --- Model~B, KPZ, Burgers
--- represent those vertices through an \term{operator IR} (intermediate
representation) of binding nodes for the Laplacian $\Lap$, the time derivative
$\Dt$, and the spatial gradient $\partial_x$. The per-vertex \emph{form-factor
table} that records which momentum form factor each interaction node carries is
populated as a \emph{side effect of evaluating the action lambda}: it is built
inside \code{model['action'](ns)}, which runs during
\code{FieldTheory.expand()}, \emph{not} during \code{\_build\_namespace}. A
cache load deliberately skips \code{expand()} --- that is the whole point --- so
the loaded namespace would carry \emph{no} table, and every derivative vertex
would silently collapse to form factor $1$. The docstring at
\file{pipeline/\_expand\_cache.py:263--264} spells out the failure mode: a
Model-B$\oplus$KPZ one-loop variance would ``load as the bare $\tilde\varphi
\varphi^2$ number instead of the correct composite\,+\,per-leg one.''

The cure is to re-run the action lambda on load. It is cheap (it builds the
symbolic action and lowers the IR --- no Taylor expansion) and reproduces the
table bit-identically. The helper is \code{\_reconstruct\_operator\_ir\_table}
(\file{pipeline/\_expand\_cache.py:253--282}):

\begin{lstlisting}
ns = getattr(ft, '_ns', None)
if ns is None or not getattr(ns, '_operator_ir', False):
    return
if hasattr(ns, '_operator_ir_vertex_terms'):
    return
try:
    ft.model['action'](ns)
except Exception:
    # Leave the table absent; the signature check below (or a
    # downstream fresh expand) catches the resulting inconsistency.
    pass
\end{lstlisting}

It is \code{hasattr}-guarded (a no-op when the table is already present), and it
returns immediately for theories that do not use the operator IR. On an
exception it swallows the error (\code{pass}) and leaves the table absent,
trusting the signature check to catch the inconsistency.

\subsection{The form-factor signature}

Because the slug cannot encode this per-vertex state, the bundle stores a
\emph{fingerprint} of it, and the load path \emph{validates} that fingerprint.
This is \code{vertex\_form\_factor\_signature}
(\file{pipeline/\_expand\_cache.py:208--250}). For a theory that does not use
the operator IR it returns \code{None} --- so every pre-existing plain or
temporal theory keeps a \code{None}-versus-\code{None} match and loads
unchanged. For an operator-IR theory it returns
\code{'operator\_ir:'} followed by the first sixteen hex digits of a SHA-256
hash, built from each vertex term's four identifying pieces: its \code{mode},
its physical-leg count \code{n\_phys}, its operator \code{chain}, and its
coupling-mixing \code{weight}.

\begin{lstlisting}
canon: list = []
for t in (terms or []):
    n_phys = t.get('n_phys')
    canon.append([
        str(t.get('mode')),
        int(n_phys) if n_phys is not None else None,
        _canon_chain(t.get('chain')),
        str(t.get('weight')),
    ])
# Order-independent: the table-build order is not contractual.
canon.sort(key=lambda e: json.dumps(e, sort_keys=True))
blob = json.dumps(canon, sort_keys=True)
return 'operator_ir:' + hashlib.sha256(blob.encode('utf-8')).hexdigest()[:16]
\end{lstlisting}

Three implementation choices deserve a word. First, every piece is stringified
(via \code{\_canon\_chain} for the operator chain, \code{str()} for mode and
weight) so it round-trips identically through pickle and across Sage sessions.
Second, the list of terms is \emph{sorted} before hashing --- using
\code{json.dumps(e, sort\_keys=True)} as the sort key --- because the
table-build order is not contractual, so the signature must be
order-independent. Third, \code{hashlib} is the Python standard-library hashing
module; \code{sha256(blob).hexdigest()[:16]} produces a short, stable
fingerprint. Adding the KPZ $(\partial_x\varphi)^2$ vertex to a Model-B-only
theory, dropping it, or re-weighting the mix all change this signature, so a
stale on-disk bundle whose signature does not match the freshly built model is
\emph{rejected} by \code{load\_expand}.

\begin{note}
The docstring at \file{pipeline/\_expand\_cache.py:229--233} records a stability
guarantee: the \code{weight} is serialised via \code{str()} of an
already-\code{simplify\_full}'d \code{SR} expression. A purely cosmetic
re-ordering of that expression across Sage versions would, at worst, force one
extra fresh \code{expand()} (and then a re-save with the new signature). It can
never cause a \emph{wrong} load --- the failure mode is conservative.
\end{note}

\subsection{The structural-integrity checks in \texttt{load\_expand}}

\code{load\_expand} (\file{pipeline/\_expand\_cache.py:351--482}) ties it all
together. Its contract is strict: it returns \code{True} on success (after which
\code{ft} is equivalent to a fresh \code{expand()} at \code{target\_order}), and
\code{False} on \emph{any} problem --- no usable cache, a load exception, or a
failed sanity check --- leaving \code{ft} unchanged so the caller can fall back
to a fresh expand. It runs a sequence of guards:

\begin{enumerate}
  \item \textbf{Discover the file.} If \code{cached\_order} was not supplied,
        re-discover it via \code{find\_best\_cached\_order}; if still
        \code{None}, return \code{False}. Then \code{sage\_load} the bundle
        inside a \code{try}/\code{except} (\file{:379--384}).
  \item \textbf{Ring-name check} (\file{:388--394}). The freshly built ring's
        generator names must equal \code{bundle['ring\_var\_names']},
        \emph{order included}. A mismatch means a model edit changed the field
        list and invalidated the cache; return \code{False}.
  \item \textbf{\code{n\_tilde} check} (\file{:396--399}).
        \code{bundle['n\_tilde']} must match \code{ft.\_n\_tilde}.
  \item \textbf{Form-factor signature check} (\file{:411--418}). Reconstruct the
        live table, then require \code{vertex\_form\_factor\_signature(ft.\_ns)}
        to equal \code{bundle['vertex\_signature']}. This catches both a changed
        model that re-used the slug \emph{and} any pre-signature bundle for a
        derivative-vertex theory.
  \item \textbf{Downgrade filter} (\file{:421--423}). If
        \code{cached\_order > target\_order}, filter the bigrades via
        \code{downgrade\_by\_tp\_dict}.
  \item \textbf{Rehydrate} (\file{:425--435}). Set
        \code{ft.\_by\_tp = \_by\_tp\_from\_dict\_form(ft.\_R, by\_tp\_dict)};
        rebuild \code{\_S\_raw} as the \emph{sum} of the (filtered) \code{by\_tp}
        polynomials; and rebuild \code{\_mf\_sector\_raw} from its dict form.
\end{enumerate}

There is one more reconstruction, and missing it would be catastrophic for
colored-noise theories. A fresh \code{expand()} calls
\code{\_build\_cumulant\_action}, which has two effects: it bakes a symbolic
noise-cumulant term into the action (captured by the cached \code{\_by\_tp}),
\emph{and} it populates \code{ns.\_cumulant\_kernels} with kernel-function
closures that the downstream type extractor reads to promote a plain noise
source into a colored one. Those closures do not round-trip through pickle, so
the load path re-executes \code{\_build\_cumulant\_action} purely for that side
effect (\file{:464--472}):

\begin{lstlisting}
try:
    from msrjd.core.field_theory import _build_cumulant_action
    _build_cumulant_action(ft._ns, model)
except Exception as e:
    if verbose:
        print(f'[expand-cache] _cumulant_kernels rebuild failed ...')
    return False
\end{lstlisting}

A failure here returns \code{False} --- forcing a fresh expand --- because a
missing \code{\_cumulant\_kernels} would silently collapse every colored-noise
diagram to zero. This, with the form-factor signature, is the protection that
keeps the slug's under-determination from producing wrong numbers.

\code{load\_expand} has one precondition the caller must honour:
\code{ft.\_ns}, \code{ft.\_R}, and \code{ft.\_n\_tilde} must already be
populated, because the dict-form rehydration needs a live ring to coerce into.
The helper \code{prepare\_for\_load} (\file{pipeline/\_expand\_cache.py:300--319})
does exactly that --- it runs \code{ft.\_build\_namespace()} (but \emph{not}
\code{expand()}) to set the namespace, ring, and \code{n\_tilde}, then calls
\code{\_reconstruct\_operator\_ir\_table}. Both production callers,
\code{precompute} and \code{compute\_cumulants}, call \code{prepare\_for\_load}
immediately before \code{load\_expand}.

\begin{gotcha}
The expand cache's on-disk slug under-determines the bundle. The filename is
only \code{model['name']} plus the taylor order; it captures neither the
operator-IR form-factor table nor the \code{\_cumulant\_kernels} closures. Both
are reconstructed on load, and the form-factor one is validated by the
signature check. If either reconstruction silently no-ops, a derivative-vertex
or colored-noise theory loads a \emph{wrong} (collapsed) result. Three separate
docstring blocks in \file{\_expand\_cache.py} are devoted to this corner ---
that is how fragile it is.
\end{gotcha}

\section{\texttt{precompute}: the one-time structural pass}
\label{sec:precompute}

Recall that bigrade order~2 already contains the $(0,0)/(1,0)/(0,1)$ mean-field
sectors \emph{and} the $(1,1)$ propagator kernel --- everything the structural
validation and propagator construction need, regardless of how high the
downstream cumulant calculation reaches. \code{precompute}
(\file{pipeline/\_precompute.py:55}) exploits this: it is a cheap (a few
seconds for any theory) one-time pass that expands at order~2, runs the sanity
check, solves the mean-field saddle, and builds the propagator --- writing
\file{expand\_taylor2.sobj} and \file{propagator.sobj} so that all later runs
get the propagator and MF check for free.

\code{precompute(model, *, force=False, verbose=True)} returns a status
\code{dict} and proceeds in four stages:

\begin{enumerate}
  \item \textbf{Expand at order~2} (\file{\_precompute.py:106--139}). Build
        \code{ft = FieldTheory(model, taylor\_order=2)}. Unless \code{force},
        try the expand cache --- \code{find\_best\_cached\_order(model, 2)}, and
        on a hit \code{prepare\_for\_load} then \code{load\_expand}. On a miss,
        \code{ft.expand()} fresh (catching exceptions into an
        \code{'EXPAND\_RAISED'} status) and \code{save\_expand} (a save failure
        is non-fatal).
  \item \textbf{Sanity check} (\file{\_precompute.py:142--152}).
        \code{ft.sanity\_check()} confirms the MF sector vanishes at the saddle;
        sets \code{mf\_check} to \code{'PASS'} or \code{'FAIL'} and returns
        early on failure.
  \item \textbf{Solve the mean field} (\file{\_precompute.py:155--164}). Build
        the parameter defaults and solve the saddle. Non-fatal on failure --- the
        propagator can still be built.
  \item \textbf{Build the propagator} (\file{\_precompute.py:167--176}).
        \code{build\_propagator(ft, model, use\_cache=True, force=force)} writes
        \file{propagator.sobj} and sets \code{propagator\_built}.
\end{enumerate}

The status dict carries \code{mf\_check} (\code{'PASS'}/\code{'FAIL'}/error
string), \code{sanity\_ok}, \code{mf\_values}, \code{taylor\_order} (always
\code{2}, since this is the pre-compute order), \code{cache\_dir},
\code{propagator\_built}, \code{wall\_seconds}, and a \code{log} list of the
status messages.

\begin{note}
The imports of \code{FieldTheory}, \code{build\_propagator}, and
\code{\_expand\_cache} sit \emph{inside} \code{precompute}
(\file{\_precompute.py:91--100}), not at module scope. This keeps importing the
module itself cheap, and lets a missing heavy dependency degrade to an
\code{'IMPORT\_FAILED'} status string rather than blowing up at import time.
\end{note}

The two internal helpers are small. \code{\_build\_fundamental\_defaults}
(\file{\_precompute.py:186}) assembles a parameter dict from each model
parameter's \code{default=} declaration, \emph{skipping} parameters marked
\code{mean\_field} (those are saddles, solved not configured) and any whose
default is absent or \code{None}. \code{\_solve\_mf\_at\_saddle}
(\file{\_precompute.py:200}) runs the DAE solver if the model declares
\code{equations}, otherwise the legacy iteration solver, and returns the saddle
values.

\section{The enumeration cache}
\label{sec:enumcache}

The second cache wraps the four-stage diagram pipeline. It lives in
\file{pipeline/\_diagrams.py} and exposes one headline function,
\code{enumerate\_unique\_diagrams}.

\subsection{The four stages it wraps}

For each loop order $\ell$, the pipeline produces the list of
\emph{non-isomorphic typed causal diagrams} with $k$ external legs, plus a
per-diagram multiplicity. The four stages (imported at
\file{\_diagrams.py:28--34} and run at \file{\_diagrams.py:176--186}):

\begin{enumerate}
  \item \textbf{Prediagrams.} \code{enumerate\_prediagrams\_all(k=k, ell=ell)}
        generates all undirected topologies --- trees, then topologies, then
        prediagrams --- with $k$ external legs and $\ell$ loops, \emph{before}
        any field assignment. (Only the prediagrams of its four-tuple return are
        kept.)
  \item \textbf{Typing.} \code{enumerate\_all\_typed} assigns each vertex an
        interaction type from the theory's vertex list (\code{vtypes}) or a
        source type (\code{stypes}), and each internal edge a propagator type
        consistent with the zero/nonzero pattern of the propagator matrix
        \code{G\_ft}. Many prediagrams admit several typings; many admit none.
  \item \textbf{Causality.} \code{filter\_causal} drops diagrams that violate
        the retarded (It\^o/causal) structure of the \msrjd{} propagator --- for
        instance an acausal response-to-response loop. It returns
        \code{(causal, n\_discarded, \_)}.
  \item \textbf{Deduplication.} \code{deduplicate\_with\_multiplicities} merges
        diagrams that are isomorphic as coloured graphs, returning the surviving
        representatives and the size of each equivalence class.
\end{enumerate}

\subsection{The equivalence relation, and what the multiplicity is for}

The deduplication step is where the genuinely subtle physics lives, and it pays
to be precise. Two typed diagrams are treated as duplicates exactly when they
have the same \term{diagram signature}: the canonical form of their coloured
incidence digraph, with leaves coloured by field only
(\file{msrjd/diagrams/symmetry.py:467} onward, called with
\code{fix\_external=False}). The canonical form is computed by Sage's
\code{DiGraph.canonical\_label(\dots, certificate=True)}, which delegates to
\code{nauty}.

\begin{note}
\code{nauty} (``No AUTomorphisms, Yes?'') is a C library by Brendan McKay for
computing graph automorphism groups and canonical labellings. Given a coloured
graph it returns a canonical form such that two graphs have the \emph{same}
canonical form if and only if they are isomorphic. SageMath embeds \code{nauty}
and exposes it through \code{Graph}/\code{DiGraph} methods like
\code{canonical\_label}. This caching subsystem never imports \code{nauty}
directly --- it inherits the dependency transitively through the
\code{deduplicate\_with\_multiplicities} call.
\end{note}

Now the subtle part. The \term{symmetry factor} $\Sym(\Gamma)$ of a diagram ---
its Feynman-rule combinatorial weight --- is carried entirely by a separate
function, \code{combinatorial\_factor}, as the orbit--stabiliser count
\[
  \Sym(\Gamma)
  \;=\; \frac{\prod_v n_{\mathrm{leg}}!}{\big|\Aut_{\text{fixed-ext}}(\Gamma)\big|},
\]
computed (\file{msrjd/diagrams/symmetry.py:438--439}) as
\code{\_wick\_leg\_factor(td)} divided by the automorphism order with external
legs held fixed. The dedup \emph{multiplicity} --- the size of each equivalence
class --- is therefore \emph{diagnostic only}. Multiplying by it would
\emph{double-count}, because the merged copies are isomorphic and their pairings
are already inside $\Sym(\Gamma)$. The multiplicity is retained in the cache
purely for diagnostics and historical compatibility (it once compensated for an
\emph{incomplete} signature --- see the Gotchas).

\subsection{The headline function and its cache key}

\code{enumerate\_unique\_diagrams} (\file{\_diagrams.py:54}) takes the expanded
\code{ft} (only \code{ft.taylor\_order} is read here, for the key), the
\code{model} (for its name), $k$ and \code{max\_ell}, the
\code{external\_fields}, the propagator \code{G\_ft}, the field-index maps
\code{resp\_idx}/\code{phys\_idx}, and the type lists \code{vtypes}/\code{stypes}.
It returns the triple \code{(unique\_by\_ell, multiplicity\_by\_ell,
all\_unique)}: a dict from $\ell$ to the surviving diagrams, a parallel dict of
their equivalence-class sizes, and a flat $\ell$-ordered concatenation.

The cache directory is per-theory, computed by \code{\_model\_cache\_dir}
(\file{\_diagrams.py:42}), which uses the very same slug rule as the expand
cache so all artefacts share one directory. The cache key, though, is finer
than the expand cache's --- it must distinguish $k$, $\ell$, external fields,
\emph{and} taylor order. The stage name (\file{\_diagrams.py:150}) is:

\begin{lstlisting}
stage_name = f'unique_typed_mult_v3_{ext_tag}_taylor{ft.taylor_order}'
\end{lstlisting}

with $k$ and $\ell$ appended by \code{PipelineCache} as \code{\_k<k>\_l<l>}.
Read the pieces:

\begin{description}
  \item[\code{unique\_typed}] the stage --- deduplicated typed diagrams.
  \item[\code{\_mult\_}] records that the file carries multiplicities.
  \item[\code{\_v3\_}] the cache-\emph{version} tag (see below).
  \item[\code{\_<ext\_tag>}] the external-field tag from
        \code{\_ext\_fields\_tag} (\file{\_diagrams.py:37}), e.g.
        \code{[('v',1),('v',2)]} becomes \code{v1\_v2}. This is what makes
        different external-field permutations land in distinct cache files
        automatically.
  \item[\code{\_taylor<N>}] the taylor-order dependence, so sibling files for
        different orders coexist in one directory.
\end{description}

The loop body (\file{\_diagrams.py:156--209}) is the manual memoise pattern. On
each $\ell$: if \code{use\_cache and cache.exists(\dots)}, try to \code{load}
the \code{\{'unique', 'multiplicities'\}} dict, populate the return structures,
and \code{continue}; any exception prints a \code{cache load failed \dots
rebuilding} message and falls through to recompute. Otherwise build the four
stages, and if \code{use\_cache} \code{save} the result --- wrapped in
\code{try}/\code{except} so a save failure only warns.

\subsection{Version tags are load-bearing}

The \code{\_v3\_} in the stage name is not decoration. It is the \emph{only}
invalidation mechanism for the enumeration cache: old files are never read
because their filename stem differs. The long comment block at
\file{\_diagrams.py:123--149} records the history. The tag was bumped
\code{v1}$\to$\code{v2} (2026-05-26) when the source-type extractor began
promoting cross-/auto-cumulant monomials to a colored \code{NoiseSourceType}; a
\code{v1} cache would encode the plain type and silently collapse every
colored-noise diagram to the white-noise answer. It was bumped
\code{v2}$\to$\code{v3} (2026-06-10) when \code{diagram\_signature} became a
\emph{complete} isomorphism invariant; a \code{v2} cache had been deduplicated
with the old incomplete signature, which collided non-isomorphic diagram classes
at $k\ge 3$ and silently dropped their integrals.

\begin{gotcha}
If you change the dedup or typing semantics, you \emph{must} bump the version
string in \code{stage\_name}. A stale cache with the old stem will load and
silently produce \emph{wrong numbers}. The \code{v2}$\to$\code{v3} bump is the
concrete cautionary tale: with the old incomplete signature, a particular
three-point cumulant came out as $-\tfrac{68}{3}\,a^3$ instead of the correct
$-32\,a^3$, because four of the eleven one-loop classes at $k=3$ had been merged
into the wrong representatives.
\end{gotcha}

\section{The bridge between the two caches: the taylor-order formula}

The two caches are linked by one formula. By default
(\file{pipeline/compute.py:279}):

\begin{lstlisting}
taylor_order = max(k + 2 * max_ell, 2)
\end{lstlisting}

The intuition: a connected diagram with $k$ external legs and $\ell$ loops needs
interaction vertices whose total leg-count covers the $k$ external legs plus the
internal propagator legs. Each loop adds (heuristically) two leg-degrees of
internal structure, giving the $2\,\texttt{max\_ell}$ term; the floor of $2$
guarantees the propagator sector $(1,1)$ is always present, even at
$k=2,\ \texttt{max\_ell}=0$. This is the taylor order needed so that
\emph{every} diagram at $(k,\ell)$ has its vertices in scope.

This formula is the bridge: it tells you which \code{taylor\_order} the expand
cache must serve to satisfy a given diagram request, and it determines the
\code{\_taylor<N>} suffix in the enumeration-cache filename. A historical floor
of $4$ was dropped to $2$ once the cache layout could side different orders
cleanly --- worth $\sim$90 minutes on a heavy Bernoulli theory at
$k=2,\ \texttt{max\_ell}=0$, because that run now needs only the cheap order-2
expansion.

\section{Data-flow walk-through}

Put the pieces together with one full
\code{compute\_cumulants(model, k=2, max\_ell=1, \dots)} run, end to end through
both caches.

\begin{enumerate}
  \item \textbf{Taylor order chosen} (\file{compute.py:279}):
        $\texttt{taylor\_order} = \max(2 + 2\cdot 1,\, 2) = 4$.
  \item \textbf{Expand cache} (\file{compute.py:319--332}). Build
        \code{ft = FieldTheory(model, taylor\_order=4)}, then
        \code{cached\_order = find\_best\_cached\_order(model, 4)} --- which is
        \code{4} if an order-4 file exists, \code{6} if only a higher order is
        cached (then downgrade-filtered), or \code{None} on a miss. On a hit,
        \code{prepare\_for\_load(ft)} then
        \code{load\_expand(model, ft, target\_order=4, cached\_order=\dots)}
        populates \code{ft.\_by\_tp}. On a miss, \code{ft.expand()} then
        \code{save\_expand(model, ft)} writes \file{expand\_taylor4.sobj}.
  \item \textbf{Downstream of expand.} The vertex/source type extractors, the
        propagator builder, and the field-index-map builder all read
        \code{ft.\_by\_tp} to produce \code{vtypes}, \code{stypes},
        \code{G\_ft}, and \code{resp\_idx}/\code{phys\_idx}.
  \item \textbf{Enumeration cache} (\file{compute.py:649}).
        \code{enumerate\_unique\_diagrams(\dots, use\_cache=True)} with
        \code{external\_fields=[('v',1),('v',2)]} forms
        \code{stage\_name = 'unique\_typed\_mult\_v3\_v1\_v2\_taylor4'}. For
        $\ell=0$ it consults
        \file{unique\_typed\_mult\_v3\_v1\_v2\_taylor4\_k2\_l0.sobj} (hit:
        load; miss: build four stages, save); for $\ell=1$ the matching
        \code{\_k2\_l1} file. It returns
        \code{unique\_by\_ell=\{0:[\dots], 1:[\dots]\}}, the parallel
        multiplicities, and the flat \code{all\_unique}.
  \item \textbf{Consumers} (\file{compute.py:680+}). Walk \code{unique\_by\_ell}
        per $\ell$, classify each diagram's coefficient factors, group by
        kernel, and integrate.
\end{enumerate}

The \code{precompute} flow is the same but truncated: it exercises only the
expand cache at order~2 and the propagator cache, writing
\file{expand\_taylor2.sobj} and \file{propagator.sobj}. A later
\code{compute\_cumulants} at order~2 then hits both immediately; at order~4 it
still pays the order-4 \code{expand()} once.

\begin{note}
The \term{spatial bridge} (\file{msrjd/integration/spatial/pipeline\_bridge.py:242})
calls \code{enumerate\_unique\_diagrams} with \code{use\_cache=False}: it
re-enumerates every time, because the diagram topology is cheap relative to the
spatial integral and the bridge does not want to manage cache keys. Note that
\code{use\_cache=False} also means \emph{never write} --- such a run does not
pollute the cache, but gets no persistence either.
\end{note}

\section{Gotchas and caveats}

A consolidated list of the sharp edges, most of which we have met in passing.

\begin{gotcha}
\textbf{Fork-based \code{multiprocessing} is dangerous on macOS notebooks.}
\code{enumerate\_unique\_diagrams} exposes \code{parallel=True}, which fans the
type-assignment stage across a \emph{fork-based} pool. Forking a Jupyter kernel
after the GUI/BLAS libraries have initialised can crash the kernel \emph{and}
the OS; the analogous temporal path was guarded for exactly this reason, and the
spatial path was made serial-only. The default here is \code{parallel=False}. Do
not flip \code{parallel=True} from a notebook on macOS --- the cache wrapper
does not re-apply the notebook-safety guard; it trusts the caller. (Cache hits
skip the parallel path entirely, since they never reach the typing stage.)
\end{gotcha}

\begin{gotcha}
\textbf{A model edit that does not change the name silently re-uses the cache.}
The slug derives from \code{model['name']} only. Editing the field list or a
vertex without bumping the name leaves a stale \file{expand\_taylor<N>.sobj} and
stale enumeration files. The partial defences are the \code{ring\_var\_names}
and \code{n\_tilde} structural checks (which catch field-list changes that alter
the ring) and the form-factor signature (which catches operator-IR vertex
changes) --- but a change that alters a coupling \emph{value} baked into a
coefficient while keeping the same ring would \emph{not} be caught. The
docstring advice (\file{\_diagrams.py:18--22}): bump \code{model['name']} or
delete the cache directory on model edits.
\end{gotcha}

\begin{gotcha}
\textbf{\code{\_model\_cache\_dir} ignores its \code{taylor\_order} argument.}
The directory is per-theory; the parameter is in the signature but never read ---
a vestige of an older per-taylor-order directory layout. Harmless, but confusing
the first time you read the call.
\end{gotcha}

\begin{gotcha}
\textbf{Graceful degradation differs slightly between the two caches.} On a load
problem the expand cache returns \code{False} (a clean miss, triggering a fresh
expand), whereas the enumeration cache \emph{rebuilds the single $\ell$ slot}
(\file{\_diagrams.py:170--173}). Both degrade gracefully and never raise out of
a cache miss --- but the granularity is per-bundle for one and per-loop-order
for the other.
\end{gotcha}

In short: the cache key is structural, the expand cache trades disk for a
ninety-minute expansion, the enumeration cache trades disk for a combinatorial
blow-up, the superset theorem makes downgrade exact, and the version string and
form-factor signature are the two guards standing between a stale file and a
wrong number. Mind those two, and the rest takes care of itself.
